\documentclass{article}
\PassOptionsToPackage{numbers,sort&compress}{natbib}
\usepackage[preprint]{neurips_2026}
\usepackage[T1]{fontenc}
\usepackage[utf8]{inputenc}
\usepackage{amsmath,amssymb}
\usepackage{hyperref}
\hypersetup{hidelinks}
\title{Neural Memory and Interval Calibration}
\author{Chengxuan \quad DSAA 2012, HKUST(GZ)}
\begin{document}
\maketitle

\section{Motivation and research question}
Financial prediction intervals should adapt to rising volatility without staying wide after a temporary shock. Neural control, score--threshold adaptation, and retrospective updates already address adaptation \citep{ncc,adapt,retro}. I will test whether \textbf{a neural forecaster's scale response changes the effect of calibration memory at shock entry and exit}. The interaction may depend on architecture and event duration. Longer context does not necessarily imply longer effective memory.

\section{Method and evaluation}
Probabilistic GRU and causal TCN models trained from scratch will produce scales $s_t$. With a fixed zero center in controlled simulations, the interval half-width is
\[
 h_t=s_t\,Q_{1-\alpha}\bigl(\{|y_i|/s_i\}_{i=t-W}^{t-1}\bigr),\qquad \alpha=0.1.
\]
Across five matched training seeds, I will cross short/long contexts with short/long score windows, using shared weights within each fit and paired event/no-event trajectories. A frozen grid varies pulse duration, magnitude, tails, and latent stochastic volatility; true scales are evaluation oracles only. Separately varying current and archived scales will distinguish their effects. Controls include EWMA/fitted GARCH and rolling, adaptive, and strongly adaptive calibration \citep{saocp}.

Public ECB reference rates for four EUR-cross currencies provide an application check: training through 2018, validation in 2019--2020, and exploratory evaluation from 2021. Metrics include interval score, width, simulated conditional coverage, and entry/recovery contrasts. Training seeds and paired paths quantify simulation variation; date-block comparisons retain cross-asset dependence. Predictions remain immutable; tuning excludes test data and event labels.

\section{Expected contribution and feasibility}
The intended contribution is a falsifiable account of \emph{when forecaster response matters for calibration}. Matched models run locally; architecture-dependent and null outcomes remain possible. Validation lag sensitivity and a finite-memory erasure check help distinguish learned response from implementation errors. Reproducible failures may motivate an adaptation module tested on a new holdout; null results will narrow the claim. November covers robustness and ablations, followed by the report and presentation before December 11.

\begin{thebibliography}{4}
\small
\bibitem[Li and Rodr\'{i}guez(2025)]{ncc}
R. Li and A. Rodr\'{i}guez. Neural Conformal Control for Time Series Forecasting. \emph{AAAI}, 2025.
\bibitem[Tumu et~al.(2026)]{adapt}
R. Tumu, A. Singh, and R. Mangharam. AdaptNC: Adaptive Nonconformity Scores for Conformal Prediction under Distribution Shift. \emph{arXiv:2602.01629v2}, 2026.
\bibitem[Jun and Ohn(2026)]{retro}
J. Jun and I. Ohn. Online conformal inference with retrospective adjustment for faster adaptation to distribution shift. \emph{Pattern Recognition}, 180:114406, 2026.
\bibitem[Bhatnagar et~al.(2023)]{saocp}
A. Bhatnagar, H. Wang, C. Xiong, and Y. Bai. Improved Online Conformal Prediction via Strongly Adaptive Online Learning. \emph{ICML}, 2023.
\end{thebibliography}
\end{document}
